\section{A prior-uniform multistep law for observation algebras}
\label{sec:algebra-multistep}
Likelihood multiplication in a fixed space gives a saturated geometry.
We prove its memory law directly and retain the transfer proof as an
alternative in Appendix~\ref{sec:algebra-transfer-alternative}. We give an experiment-level criterion: neither the
global cover nor the acquired flag is an additional hypothesis. In this
class the covariance spectrum controls every checkpoint, including
simultaneous degenerations of the prior and the detector. The parameter
space and geometric argument differ from those of the monomial theorem.

\subsection{Experiments and the covariance profile}
Let $X$ be a nonempty compact Hausdorff space, $d\ge1$, and
$f=(f_1,\ldots,f_d)^t\in C(X)^d$, with $|f_i|\le1$.
Assume that for $1\le i,j\le d$ there are specified real coefficients
\begin{equation}\label{eq:algebra-multiplication}
 f_if_j=a_{ij}+\sum_{k=1}^d b_{ijk}f_k,
 \qquad |a_{ij}|,|b_{ijk}|\le B.
\end{equation}
The functions may be dependent, including dependence modulo constants.
Only the existence of the displayed uniformly bounded representations is
required. In a parameter family these representations are part of the
known calibration. Put $W=\operatorname{span}\{1,f_1,\ldots,f_d\}$.

Fix $N\ge2$ and $0<\delta\le1/(4d)$. At each acquisition trial a
command $u\in U=[-\delta,\delta]^d$ gives a binary report $y\in\{-1,1\}$
with likelihood
\begin{equation}\label{eq:algebra-likelihood}
 L_y(u,x)=\tfrac12(1+y\,u^tf(x)).
\end{equation}
Thus $3/8\le L_y\le5/8$. Conditional on the one latent variable,
trials are independent. The exploration commands are independent and
uniform on $U$. Let $\mu$ be any Borel probability on $X$, not necessarily
of full support, and set
\begin{equation}\label{eq:algebra-covariance}
 m=\mu f,\qquad \Sigma=\mu[(f-m)(f-m)^t],\qquad
 \lambda_1\ge\cdots\ge\lambda_d\ge0
\end{equation}
for its covariance and eigenvalues. The prior is known and fixed during
a run. All uniform assertions below allow it to vary between runs.

At checkpoint $n$, $1\le n<N$, put $s=N-n$. There are $d+1$ possible
future queries, chosen uniformly and independently \emph{after} the
label is retained. Query $0$ uses the zero command at all $s$ trials.
Query $i$, $1\le i\le d$, uses $\delta e_i$ at the first trial and
zero at the others. Its outcome is the indicator of the all-positive
report word; only the chosen query is executed, consuming $s$ trials.
The query success functions are
\begin{equation}\label{eq:algebra-query}
 H_{s,0}=2^{-s},\qquad H_{s,i}=2^{-s}(1+\delta f_i).
\end{equation}
They span $W$ together with their constant component. In fact they span
the space of all $s$-trial likelihood products: closure proves one
inclusion, and \eqref{eq:algebra-query} proves the other. Adaptive
future words give the same space, since commands on a fixed word are
fixed. Squared prediction regret is evaluated for this physical menu,
not for covariance-whitened coordinates.

For a feasible prefix $h$, write $v_n(h)$ for its posterior mean and
$p_{n,i}(h)=\mu(H_{N-n,i}\mid h)$ for its physical query probabilities.
An encoder--decoder pair with at most $M$ labels produces, through its
label, predictions $a_{n,i}(h;\omega)\in[0,1]$. Here $\omega$ denotes
coding randomness independent of the exploration experiment. A public
seed may specify the codebook but cannot carry acquired information;
private coins have the resource convention of
Definition~\ref{def:finite-state}. Set
\begin{equation}\label{eq:algebra-local-loss}
 \ell_n(h;\omega)=\frac1{d+1}\sum_{i=0}^d
                  |a_{n,i}(h;\omega)-p_{n,i}(h)|^2.
\end{equation}
The expectation over the future binary outcome has already been taken:
\eqref{eq:algebra-local-loss} is its squared-loss excess.

\begin{definition}[Checkpoint and causal risk criteria]
\label{def:algebra-risks}
Let $\mathcal H_n$ be the set of feasible prefixes and $\mathsf P_n$
their law under the independent uniform commands and actual reports.
Let $\mathcal A_{M,n}$ be all measurable checkpoint encoder--decoder
pairs using at most $M$ labels. Define
\begin{align}
 R^{\rm av}_{M;n}
  &=\inf_{A\in\mathcal A_{M,n}}
                 \E_{h\sim\mathsf P_n,\omega}\ell_n^A(h;\omega),
       \label{eq:algebra-risk-average}\\
 R^{\rm max}_{M;n}
  &=\inf_{A\in\mathcal A_{M,n}}\sup_{h\in\mathcal H_n}
                 \E_\omega\ell_n^A(h;\omega).
       \label{eq:algebra-risk-worst}
\end{align}
For $\boldsymbol M=(M_1,\ldots,M_{N-1})$, let
$\mathcal C_{\boldsymbol M}$ be the class of single causal filters with
these stage budgets, including their decoders. Put
\begin{align}
 \mathfrak R^{\rm av}_{\boldsymbol M}
  &=\inf_{F\in\mathcal C_{\boldsymbol M}}\max_{1\le n<N}
                  \E_{h\sim\mathsf P_n,\omega}\ell_n^F(h;\omega),
       \label{eq:algebra-risk-causal-average}\\
 \mathfrak R^{\rm max}_{\boldsymbol M}
  &=\inf_{F\in\mathcal C_{\boldsymbol M}}\max_{1\le n<N}
       \sup_{h\in\mathcal H_n}\E_\omega\ell_n^F(h;\omega).
       \label{eq:algebra-risk-causal-worst}
\end{align}
All $\mathsf P_n$ are marginals of one exploration law. The infimum is
outside the checkpoint maximum: different checkpoint optima are not
identified with one causal filter. In the worst-history criteria coding
randomness is averaged before the supremum. Only one selected checkpoint
and its query are executed in a run.
\end{definition}

Define
\begin{equation}\label{eq:algebra-profile}
 V_\ell(\Sigma)=\prod_{i=1}^\ell\lambda_i
       =\|\mathop{\bigwedge}\nolimits^\ell\Sigma\|_{\rm op},
 \qquad
 \mathcal E_\Sigma(M)=\max_{1\le\ell\le d}
                   (V_\ell(\Sigma)/M)^{2/\ell}.
\end{equation}
A zero product contributes zero.

The geometry is already saturated at the first positive checkpoint.
Writing $k=\rank\Sigma$, Proposition~\ref{prop:algebra-saturation}
proves, for fixed $r,R>0$ independent of the prior,
\[
 m+\Sigma(rB_d)\ \subset\ S_n\ \subset\ m+\Sigma(RB_d),
                    \qquad \dim S_n=k\quad(1\le n<N).
\]
Moreover, $W$ is the full algebra on a finite clopen observable
partition, by Proposition~\ref{prop:finite-observable-quotient}.
If $t$ of its parts have positive prior mass, then $k=t-1$.
The parameterization may become redundant and partition parts may merge;
no inverse of their feature separation is used. Thus this class has no
additional history-length acquisition cap below covariance rank. The
following assertion is a prior-uniform causal classification of this
saturated case, rather than a claim of increasing observable dimension.

\begin{theorem}[Prior-uniform algebra classification]
\label{thm:algebra-multistep}
Under \eqref{eq:algebra-multiplication}, there are constants $c,C,C_N>0$
depending only on $d,B,\delta,N$ such that for every $\mu$, every
integer $M\ge1$, and $1\le n<N$,
\begin{equation}\label{eq:algebra-checkpoint-law}
 c\mathcal E_\Sigma(M)\le R^{\rm av}_{M;n}
 \le R^{\rm max}_{M;n}\le C\mathcal E_\Sigma(M).
\end{equation}
For arbitrary integer stage budgets,
\begin{equation}\label{eq:algebra-causal-law}
 c\max_n\mathcal E_\Sigma(M_n)
 \le\mathfrak R^{\rm av}_{\boldsymbol M}
 \le\mathfrak R^{\rm max}_{\boldsymbol M}
 \le C_N\max_n\mathcal E_\Sigma(M_n).
\end{equation}
The last bound is attained by a deterministic causal filter. More
precisely, its retained representatives have exact-moment errors
\begin{equation}\label{eq:algebra-accumulation}
 e_n\le C\sum_{j=1}^n L^{n-j}\mathcal E_\Sigma(M_j)^{1/2},
\end{equation}
where $L$ depends only on $d,B,\delta$. The lower bounds allow independent
public or private coding randomness. All constants are independent of
the positive eigenvalues, their multiplicities, the prior, and the
particular functions subject to the displayed bounds. If $\Sigma=0$,
one label is exact at every checkpoint.
\end{theorem}

Thus the affine covariance law has a genuine multistep extension in this
class. The assertion does not follow by applying the one-step theorem
successively with newly conditioned priors: that would leave both the
common acquisition law and accumulation of coding error unproved.
We establish the whole-history and actual-mass facts for the original
experiment, then use an elementary ellipsoid cover and the raw update.

\subsection{Whole-history product coordinates}
Use the coefficient norm on $\R\oplus\R^d$. Formula
\eqref{eq:algebra-multiplication} defines a bounded bilinear rule for
multiplying represented functions:
\[
 (a,b)\star(c,z)=\left(ac+\sum_{i,j}b_i z_j a_{ij},\quad
             az+cb+\left(\sum_{i,j}b_i z_j b_{ijk}\right)_{k=1}^d\right).
\]
Its norm is bounded by a constant depending only on $d,B$.
The represented function is the pointwise product. The coefficient rule
need not be associative when the representation is redundant; fix
left-to-right bracketing throughout. The constant vector $(1,0)$ is
an exact identity for this rule.

\begin{lemma}[A bounded product chart for every history]
\label{lem:algebra-product-chart}
For every prefix $h=(u_1,y_1,\ldots,u_n,y_n)$ with $1\le n<N$, there
are coefficient polynomials $a_n(h)$ and $b_n(h)\in\R^d$ of degree at
most $n$ such that
\[
 P_h^\circ=\prod_{j=1}^n(1+y_j u_j^tf)=a_n(h)+b_n(h)^tf.
\]
The coefficients are bounded uniformly by $d,B,\delta,N$. If
$Z_n=\mu P_h^\circ=a_n+b_n^tm$ and $v_n=\mu(fP_h^\circ)/Z_n$, then
\begin{equation}\label{eq:algebra-chart}
 (3/4)^n\le Z_n\le(5/4)^n,\qquad
 v_n=m+\Sigma\theta_n(h),\qquad \theta_n(h)=b_n(h)/Z_n.
\end{equation}
In particular $\|\theta_n(h)\|\le R$ for a uniform $R$.
The full reachable set of $v_n$, not just a regular patch, is a compact
semialgebraic set of bounded format and dimension at most $d$.
\end{lemma}
\begin{proof}
Starting from $(1,0)$, apply the fixed coefficient rule to
$(1,y_j u_j)$ in chronological order. Bounded bilinearity gives the
coefficient bound inductively, and each step raises total command
degree by at most one. Evaluation on $X$ proves the product identity,
regardless of dependence among the functions. The evidence bounds follow
pointwise from $|u_j^tf|\le d\delta\le1/4$.

Since $\mu(ff^t)=\Sigma+mm^t$, integration gives
\[
 \mu(fP_h^\circ)=a_nm+(\Sigma+mm^t)b_n
                       =mZ_n+\Sigma b_n.
\]
Division by the positive evidence proves \eqref{eq:algebra-chart} and
the bound on $\theta_n$. For a fixed report word, $v_n$ is a rational
map of the compact command cube, of bounded degree, with denominator
bounded away from zero. Its image is compact. Its graph and projection
have a semialgebraic description with format bounded in terms of
$n,d$ and the degrees, by real quantifier elimination. The real
coefficients can depend on $\mu$ and $f$; no definability in those
parameters is needed. Taking the finite union over words preserves the
bound. The image is in $\R^d$, giving the asserted dimension bound.
\end{proof}

\subsection{Actual acquisition mass at every checkpoint}
The variable $\theta_n$ is a coordinate on complete histories used only
for the geometric proof. It is not an additional retained tape. At
singular $\Sigma$ it need not be determined by $v_n$.

\begin{lemma}[An unconditional product-coordinate density]
\label{lem:algebra-acquired-density}
There are $r,\gamma>0$, depending only on $d,B,\delta,N$, such that
at every checkpoint the actual law of $\theta_n$ dominates
$\gamma$ times Lebesgue measure on $rB_d$. The same independent-command
experiment supplies these bounds for all $n$. The claim includes the
probability of the selected report word and is uniform over all priors.
\end{lemma}
\begin{proof}
It suffices to use the all-positive word. Write $z=u_1$ and
$w=(u_2,\ldots,u_n)$. On this word put
$F_n(z,w)=b_n(z,w)/(a_n(z,w)+b_n(z,w)^tm)$.
At $w=0$ the coefficient identity rule gives exactly
\[
 F_n(z,0)=\frac{z}{1+m^tz},\qquad F_n(0,0)=0,
                \qquad D_zF_n(0,0)=I_d.
\]
The bounds $|m_i|\le1$, the bounded coefficient rule, the bounded
polynomial degrees and the evidence lower bound give a common bound
on all first and second command derivatives of $F_n$ on a neighborhood
of zero, for $1\le n<N$. This follows directly by differentiating
$b_n/Z_n$: its derivatives involve only the bounded polynomial
coefficients and powers of $Z_n^{-1}$. In particular, this bound does
not use $\Sigma^{-1}$ or any support condition on $\mu$.

Choose $a>0$ so that the closed ball $aB_d$ lies inside the command
cube and $\|D_zF_n-I_d\|\le1/4$ there when $w$ is sufficiently
small. Shrink a positive-volume ball $W_n$ of complementary commands
so that this derivative bound holds on $aB_d\times W_n$ and
$\|F_n(0,w)\|\le a/8$ for every $w\in W_n$.
The sizes can be chosen uniformly in all experiment parameters and
then uniformly in $n$, since there are only finitely many checkpoints.
For $n=1$, $W_1$ is the zero-dimensional singleton, of volume one.

For fixed $w\in W_n$ and $t\in(a/2)B_d$, the map
\[
 z\longmapsto t-F_n(0,w)-\bigl(F_n(z,w)-F_n(0,w)-z\bigr)
\]
is a contraction of the closed ball $aB_d$ into itself: its Lipschitz
constant is at most $1/4$, and its norm there is at most
$a/2+a/8+a/4<a$. Its unique fixed point solves $F_n(z,w)=t$.
Moreover $F_n(\cdot,w)$ is injective on $aB_d$, because its difference
from the identity is $1/4$-Lipschitz. Its derivative has singular
values between $3/4$ and $5/4$. The inverse function theorem therefore
makes it a diffeomorphism onto its image, and
$|\det D_zF_n|\le(5/4)^d$.

The joint density of all commands and this report word, before any
conditioning, is
\[
 (2\delta)^{-nd}\,2^{-n}Z_n(z,w)
                         \ge (2\delta)^{-nd}(3/8)^n.
\]
For a Borel set $A\subset(a/2)B_d$, change variables in $z$ for each
$w\in W_n$ and integrate in the complementary coordinates. The
probability that $\theta_n\in A$ and that this word occurs is at least
\[
 (2\delta)^{-nd}(3/8)^n\,\operatorname{vol}(W_n)
                                      (4/5)^d\operatorname{vol}(A).
\]
Taking the minimum of these positive constants over $n$ proves the
lemma with $r=a/2$. No fixed complementary command is assigned
positive probability: its open neighborhood has been integrated.
All prefixes are marginals of the original independent-command law;
there is no change of exploration experiment between checkpoints.
\end{proof}

\subsection{Raw updates and a direct proof}
\begin{lemma}[A collision-free algebra update]
\label{lem:algebra-update}
The raw posterior mean $v=\nu f$ updates after $(u,y)$ by
\begin{equation}\label{eq:algebra-update}
 v'_i=\frac{v_i+y\sum_j u_j\bigl(a_{ij}+\sum_k b_{ijk}v_k\bigr)}
                        {1+y\,u^tv}.
\end{equation}
This map is uniformly Lipschitz on every segment between posterior
moment states, with constant depending only on $d,B,\delta$.
It uses the current state, command and report only, and maps each
reachable state to a reachable next state. No inverse covariance or
independent basis of $W$ is required.
\end{lemma}
\begin{proof}
Insert \eqref{eq:algebra-multiplication} in Bayes' formula. The factors
$1/2$ cancel, giving \eqref{eq:algebra-update}. On the segment of
moment vectors corresponding to a convex mixture of two posteriors,
$|v_i|\le1$ and $1+y\,u^tv\ge3/4$. The numerator, denominator and
their first derivatives are bounded by $d,B,\delta$; the quotient rule
therefore gives the claimed Lipschitz bound. Integration along the
segment gives the estimate between its endpoints. Extending a history
by the admitted current command and report realizes the updated state.
This also verifies directly the hypotheses of
Lemma~\ref{lem:positive-remaining-update}.
\end{proof}

\begin{proposition}[Immediate saturation of the observable geometry]
\label{prop:algebra-saturation}
Let $S_n$ be the reachable raw posterior mean set in the algebra
experiment, and put $k=\rank\Sigma$. There exist $r,R>0$ depending
only on $d,B,\delta,N$ such that
\begin{equation}\label{eq:algebra-saturated-sandwich}
 m+\Sigma(rB_d)\subset S_n\subset m+\Sigma(RB_d),
                 \qquad 1\le n<N.
\end{equation}
In particular, $S_n$ has dimension $k$ at every positive checkpoint.
For the fixed observable partition of
Proposition~\ref{prop:finite-observable-quotient}, if exactly $t$ parts
have positive prior mass, then $k=t-1$.
\end{proposition}
\begin{proof}
The outer inclusion is Lemma~\ref{lem:algebra-product-chart}.
The inverse construction in
Lemma~\ref{lem:algebra-acquired-density} realizes every point of a
fixed ball in the $\theta_n$ coordinates by admitted commands and the
all-positive word. Shrinking its radius if necessary gives the first
inclusion, including the boundary of the smaller ball. Equivalently,
the density domination and closedness of the reachable image give the
same inclusion. The image of $\Sigma$ has dimension $k$ and
$\Sigma(rB_d)$ has nonempty relative interior in that image. The
two inclusions prove the dimension assertion, also when $k=0$.

Let $a_j\in\R^d$ be the joint feature values on the observable
partition. The rows $(1,a_j^t)$ are linearly independent: evaluation
of $W$ on all parts is onto because $W$ contains every partition
indicator. Any subset of these rows remains independent. For the
$t$ positive-mass parts, a vector $z$ satisfies $z^t\Sigma z=0$
exactly when $z^ta_j$ is the same on all $t$ parts. The affine
independence of these feature vectors shows that their difference
span has dimension $t-1$. This span is the covariance range, proving
$k=t-1$. This qualitative rank calculation uses no conditioning bound
for the indicator basis.
\end{proof}

\begin{proof}[Direct proof of Theorem~\ref{thm:algebra-multistep}]
We give the complete covering and mass argument to separate it from the
nonsaturated geometry of the monomial theorem. This is the thick-linear
mechanism of Lemma~\ref{lem:v18-thick-linear}, with the history and
causality verifications made for the present experiment.

Suppose first that $k>0$, and put
\[
 e(M)=\max_{1\le j\le k}(V_j(\Sigma)/M)^{1/j}
                         =\mathcal E_\Sigma(M)^{1/2}.
\]
Rotate the covariance ellipsoid orthogonally. Its containing rectangle
has half-widths $R\lambda_1,\ldots,R\lambda_k$. Partition each
coordinate interval into intervals of length at most
$\varepsilon/\sqrt{k}$, with at most
$1+2\sqrt{k}R\lambda_i/\varepsilon$ intervals. Every product cell
has diameter at most $\varepsilon$. Discard empty cells and choose a
point of $S_n$ in every remaining cell. This gives reachable centers
and the estimate
\begin{align}
 N(S_n,\varepsilon)
 &\le\prod_{i=1}^k
         (1+2\sqrt{k}R\lambda_i/\varepsilon)\notag\\
 &\le 1+\sum_{j=1}^k\binom{k}{j}
                (2\sqrt{k}R/\varepsilon)^jV_j(\Sigma).
       \label{eq:algebra-direct-cover}
\end{align}
Here $N$ denotes the minimum number of radius-$\varepsilon$ balls
with centers in $S_n$. Zero axes require only one coordinate value.
For $\varepsilon=Ae(M)$, each $V_j/e(M)^j\le M$.
Choose $A$, depending only on $d,R$, so that
$\sum_{j=1}^k\binom{k}{j}(2\sqrt{k}R/A)^j\le1/2$ for
all $k\le d$. Then \eqref{eq:algebra-direct-cover} is at most
$1+M/2\le M$ for every integer $M\ge2$. For $M=1$ use the
reachable state $m$, obtained by zero commands: its radius is at most
$R\lambda_1=Re(1)$. This proves the covering assertion for every
integer budget, not only along an asymptotic subsequence.

For the lower bound rotate the coefficient ball in
Lemma~\ref{lem:algebra-acquired-density} by the same orthogonal matrix.
Take $a=r/(2\sqrt d)$. Its cube $[-a,a]^d$ remains in the ball.
Projecting the dominated cube measure onto the first $j$ coordinates
and multiplying by the positive eigenvalues gives domination by a
fixed mass $\beta=\gamma(2a)^d$ times uniform probability on
\[
 Q_j=\prod_{i=1}^j[-a\lambda_i,a\lambda_i],
                       \qquad 1\le j\le k.
\]
This mass is from the original command-and-report law; the selected
word has not been conditioned away. For any $M$ centers in $\R^j$,
the fraction of $Q_j$ covered by radius-$b$ balls is at most
\[
 \frac{M\operatorname{vol}(B_j)b^j}{(2a)^jV_j(\Sigma)}.
\]
Taking $b=c_j(V_j/M)^{1/j}$ with a sufficiently small fixed $c_j$
leaves at least half of this dominated mass uncovered. The expected
squared distance is therefore at least $\beta b^2/2$.

The physical query map has the exact identity
\begin{equation}\label{eq:algebra-direct-metric}
 \frac1{d+1}\sum_{i=0}^d|p_{n,i}(v)-p_{n,i}(w)|^2
     =\frac{2^{-2(N-n)}\delta^2}{d+1}\|v-w\|^2.
\end{equation}
For arbitrary decoder predictions, subtract the constant $2^{-(N-n)}$
from the last $d$ coordinates and divide by $2^{-(N-n)}\delta$.
The result is a set of at most $M$ raw-state centers. Discarding the
constant-query error and projecting orthogonally shows that the same
lower estimate applies, with the fixed metric factor in
\eqref{eq:algebra-direct-metric}. Randomized assignment cannot improve
nearest-center distance. Average private decoder randomness conditional
on its label; condition on independent public randomness and then
integrate. The bound is unchanged. Maximizing over $j$ proves the
average checkpoint lower bound; it also bounds the worst-history
criterion. The reachable cover and
\eqref{eq:algebra-direct-metric} prove its upper bound.

For the causal assertion, fix these reachable codebooks separately for
the prescribed $M_n$, but use one filter. Starting from $m$, update the
retained representative by \eqref{eq:algebra-update} and choose its
nearest representative at the next stage, with a fixed tie rule.
Lemma~\ref{lem:algebra-update} makes the candidate reachable and gives
\[
 e_n\le Le_{n-1}+C\mathcal E_\Sigma(M_n)^{1/2},
                               \qquad e_0=0.
\]
Iteration is exactly \eqref{eq:algebra-accumulation}. Squaring and
using the finite horizon and the physical metric gives the upper bound
in \eqref{eq:algebra-causal-law}. For every single causal filter its
label at time $n$ is an admissible $M_n$-label checkpoint encoding.
The preceding lower bounds thus hold for that filter at every $n$
under the common exploration law. Taking their maximum and only then
the infimum over filters gives the causal lower bound of
Definition~\ref{def:algebra-risks}. This does not interchange an
infimum and a maximum.

If $k=0$, the product chart makes $S_n=\{m\}$ for every $n$;
its constant prediction and raw update use one label and have zero
error. Every step above has constants depending only on
$d,B,\delta,N$. No eigenvalue inverse, algebraic-format estimate, or
lower bound on a nonzero prior weight occurs.
\end{proof}



\subsection{Observable quotients and degenerating finite experiments}
The algebra assumption has a precise scope; it is not automatic for a
finite list of continuous functions.

\begin{proposition}[Finite observable quotient]
\label{prop:finite-observable-quotient}
A finite-dimensional unital real subalgebra $W\subset C(X)$ is exactly
the algebra of functions constant on a finite clopen partition of $X$.
The number of its nonempty parts is $\dim W$. In particular, the class
above allows arbitrary compact latent spaces, but its observable
quotient is finite. The partition may merge in a degenerating family;
no well-conditioned partition basis is required by the theorem.
\end{proposition}
\begin{proof}
For any $g\in W$, sufficiently many powers of $g$ are linearly
dependent in $W$. Thus a nonzero real polynomial annihilates $g$
pointwise, and $g(X)$ is a finite set of real roots. Apply this to a
finite basis of $W$. Their joint value map has finitely many values;
its nonempty fibers form a finite clopen partition, and every member
of $W$ is constant on these fibers.

Conversely, choose a joint value $a$. For each other joint value $b$,
some basis function $g_{i(a,b)}$ has different values at $a$ and $b$.
The product
\[
 \prod_{b\ne a}
 \frac{g_{i(a,b)}-g_{i(a,b)}(b)}
      {g_{i(a,b)}(a)-g_{i(a,b)}(b)}
\]
belongs to $W$ and is the indicator of the fiber indexed by $a$.
These indicators span all functions on the partition and are linearly
independent, proving both assertions. The denominators in this
qualitative argument need not be uniformly bounded away from zero.
They occur nowhere in the quantitative proof above, which uses the
given bounded multiplication representations instead.
\end{proof}

\begin{corollary}[Simultaneously vanishing acquisition and query channels]
\label{cor:algebra-attenuation}
Let $X=\{0,1,\ldots,d\}$, $\mu\{i\}=p_i\ge0$ with $\sum p_i=1$,
and $f_i=\tau_i\mathbf1_{\{i\}}$, $0\le\tau_i\le1$.
The preceding theorem holds uniformly for all $p$ and $\tau$, with
\begin{equation}\label{eq:algebra-indicator-covariance}
 \Sigma=D_\tau\bigl(\operatorname{diag}(p_1,\ldots,p_d)-p_+p_+^t\bigr)D_\tau,
 \qquad p_+=(p_1,\ldots,p_d)^t.
\end{equation}
If all $p_i\ge p_*>0$, including $p_0$, and
$\tau_{(1)}\ge\cdots\ge\tau_{(d)}$ are the sorted attenuations, then
\[
 \mathcal E_\Sigma(M)\asymp
 \max_{1\le\ell\le d}
       \left(\frac{\prod_{i=1}^\ell\tau_{(i)}^2}{M}\right)^{2/\ell}.
\]
For $d=2$ and $p_0=p_1=p_2=1/3$, the common-budget causal law is
\begin{equation}\label{eq:algebra-two-channel-law}
 \mathfrak R^{\rm av}_{M}\asymp\mathfrak R^{\rm max}_{M}
 \asymp \max\{\tau_{\max}^4M^{-2},\ \tau_1^2\tau_2^2M^{-1}\}.
\end{equation}
When $\tau_{\min}>0$, the crossover order is
$M\asymp(\tau_{\max}/\tau_{\min})^2$. At $\tau_{\min}=0$
only the one-dimensional branch remains, and at $\tau_{\max}=0$
one label is exact.
\end{corollary}
\begin{proof}
The multiplication rules are $f_i^2=\tau_i f_i$ and $f_if_j=0$
for $i\ne j$, with bounded coefficients even at zero attenuation.
Direct integration gives \eqref{eq:algebra-indicator-covariance}.
For every $z\in\R^d$, Cauchy--Schwarz gives
\[
 p_0\sum_{i=1}^d p_i\tau_i^2z_i^2
 \le z^t\Sigma z\le\sum_{i=1}^d p_i\tau_i^2z_i^2.
\]
Under the strict lower bounds on the $p_i$, this places $\Sigma$
between fixed positive multiples of $D_\tau^2$ in quadratic-form
order. The min--max characterization of ordered eigenvalues proves
the scale comparison and hence the profile formula. Without those
lower bounds, the exact covariance formula and its uniform theorem
still hold, but the attenuation-only comparison is not asserted.

In the two-channel case,
\[
 \Sigma=\frac19\begin{pmatrix}2\tau_1^2&-\tau_1\tau_2\\
                      -\tau_1\tau_2&2\tau_2^2\end{pmatrix},
 \qquad \det\Sigma=\frac{\tau_1^2\tau_2^2}{27}.
\]
Its largest eigenvalue is comparable to $\tau_{\max}^2$.
Substitution into \eqref{eq:algebra-profile} proves the two branches;
equating them gives the crossover order. Zero cases follow directly
from the covariance formula. The fourth power in the first branch
has a physical cause: the same attenuated feature participates once
in acquisition and once in the later query. It has not been normalized
away by an inverse covariance.
\end{proof}

\begin{corollary}[Operational recovery of the multistep invariant]
\label{cor:algebra-operational-inverse}
Let $\mathfrak R_M^{\rm av}$ denote the common-budget causal curve of
Theorem~\ref{thm:algebra-multistep}. For $1\le\ell\le d$,
\begin{equation}\label{eq:algebra-operational-inverse}
 \inf_{M\in\N,\ M\ge1}M(\mathfrak R_M^{\rm av})^{\ell/2}
                         \asymp V_\ell(\Sigma).
\end{equation}
Thus uniform comparison of all integer-budget causal curves is
equivalent to uniform comparison of their covariance exterior-power
products, after padding to the fixed dimension. The assertion includes
zero products and also holds for the worst-history curve.
\end{corollary}
\begin{proof}
Write $e(M)=\mathcal E_\Sigma(M)^{1/2}$. Every branch gives
$Me(M)^\ell\ge V_\ell$. If $\lambda_\ell>0$, set
$M_0=V_\ell/\lambda_\ell^\ell\ge1$. Ordering of the eigenvalues
implies $(V_j/M_0)^{1/j}\le\lambda_\ell$ for every $j$:
for $j<\ell$ divide by $\lambda_{j+1}\cdots\lambda_\ell$, and
for $j>\ell$ multiply by $\lambda_{\ell+1}\cdots\lambda_j$.
Hence $e(M_0)=\lambda_\ell$. Monotonicity and
$\lceil M_0\rceil\le2M_0$ show
\[
 V_\ell\le\inf_{M\ge1}Me(M)^\ell\le2V_\ell.
\]
If the covariance rank is $k<\ell$, then either $k=0$ and $e=0$,
or $e(M)=O(M^{-1/k})$, so the infimum is zero. The uniform forward
law now proves \eqref{eq:algebra-operational-inverse}. Comparing the
products proves comparison of the finite maximum defining the forward
curves, and the inverse proves the converse. This uses the full
unbounded sequence of integer budgets, not finitely sampled risks.
\end{proof}

The invariant has therefore been reused at the level of a complete
multistep experiment class. The covariance theorem does not extend the
exact-kernel criterion to all positive detectors, nor does it identify
finite label cardinality with total computational space. Its additional
conclusion is the prior-uniform causal law under a concrete bounded
multiplication condition. The monomial, affine, and exact-kernel
arguments retain their distinct hypotheses and full statements.
